% Generated draft: requires mathematical and source review.
\documentclass[11pt]{article}
\usepackage[margin=25mm]{geometry}
\usepackage{fontspec}
\setmainfont{Latin Modern Roman}
\usepackage{amsmath,amssymb,mathtools,unicode-math}
\setmathfont{Latin Modern Math}
\usepackage{enumitem,needspace,xurl,hyperref}
\setlength{\parindent}{0pt}
\setlength{\parskip}{5pt}
\setlength{\emergencystretch}{3em}
\newcommand{\N}{\mathbb{N}}
\newcommand{\Z}{\mathbb{Z}}
\newcommand{\Q}{\mathbb{Q}}
\newcommand{\R}{\mathbb{R}}
\newcommand{\C}{\mathbb{C}}
\newcommand{\F}{\mathbb{F}}
\newcommand{\A}{\mathbb{A}}
\newcommand{\PP}{\mathbb{P}}
\newcommand{\E}{\mathbb{E}}
\newcommand{\eps}{\varepsilon}
\newcommand{\dd}{\,\mathrm{d}}
\begin{document}
\section*{UnsolvedMath ID: 2919}
\textbf{Source problem code:} catalogue entry
\subsection{Definitions and mathematical statement}
Let \(X\) be the source's negative-definite smooth 4-manifold, with an embedded 4-ball \(B^4\), and put \(W=X\setminus\operatorname{int}B^4\). To use the classical knot invariant, take \(K\) on the distinguished boundary sphere \(\partial B^4\cong S^3\). Let \(\Sigma\subset W\) be a compact connected properly smoothly embedded orientable surface with \(\partial\Sigma=K\), and let \(g(\Sigma)\) be its genus. These are the usual genus-bound conventions; the source does not explicitly state all compactness, orientability and boundary-component hypotheses.

For an oriented \(X\), negative definite means that its real intersection form satisfies \(Q_X(v,v)<0\) for every nonzero \(v\in H_2(X;\R)\). Rank-zero forms are allowed, as the supplied \(S^4\) discussion requires. No simple-connectivity or vanishing-first-homology hypothesis is added.

The exact displayed target is
\[
2g(\Sigma)\geq|s(K)|,
\]
where \(s\) is Rasmussen's original concordance invariant, rather than one of the field-coefficient or Steenrod-refined variants mentioned separately. The source says that \(\Sigma\) is nullhomotopic. If this means the inclusion \(\Sigma\hookrightarrow W\) is nullhomotopic as an ordinary map, that condition must be retained. It is not silently replaced by
\[
[\Sigma]=0\in H_2(W,S^3;\Z),
\]
the nullhomologous relative-surface condition appearing in the supplied special-case discussion. The sense of nullhomotopy, including any relative-boundary requirement, and its intended relation to nullhomology require source review before these can be identified as one proposition.

The background and triage dated 2026-08-17 report Rasmussen's standard \(X=S^4\) result and Manolescu--Marengon--Sarkar--Willis's nullhomologous-surface results for standard negative-definite connected sums and related manifolds. The background prints connected sums of \(\mathbb{CP}^2\) without orientation bars; the supplied review explicitly identifies lost bars, so the negative-definite comparison is with the orientation-reversed projective planes, not silently with the positive-definite standard orientation.

The source reports Freedman--Gompf--Morrison--Walker's proposed route to the smooth four-dimensional Poincar\'e problem: a knot slice in a homotopy sphere but with \(s(K)\ne0\) would distinguish it from standard slicing. The precise surface hypothesis needed to apply the present inequality must still be checked; this is an approach, not a constructed counterexample.

It also reports a gauge-theoretic concordance invariant printed as \(s7\) or \(s^7\), identified by the supplied review as OCR-damaged \(s^\sharp\), with an analogous genus bound. The background says it differs from Rasmussen's invariant for known knots, including the trefoil, so that gauge-theoretic result is not the requested theorem. The field-coefficient Rasmussen variants and Lipshitz--Sarkar \(Sq^1\) refinements are separately proposed extensions, with damaged exact indexing in the source; none is silently substituted for \(s\).

The registry status is open. The supplied dated triage labels the row partially solved and leaves the exact Rasmussen inequality for arbitrary negative-definite smooth 4-manifolds open, with no general gauge-theoretic interpretation found in its check. No source, genus theorem or suggested counterexample approach is independently verified here.
\subsection{Short English statement}
Does Rasmussen's \(s\)-invariant bound the genus of the specified surface in every negative-definite smooth 4-manifold? The source reports standard special cases, but its nullhomotopic wording must be distinguished from the nullhomologous hypothesis in those results. Orientation-bar and auxiliary-invariant OCR issues also remain explicit.
\subsection{Sources}
\begin{itemize}
\item Local UnsolvedMath registry, permanent ID 2919, source code catalogue entry, ``Kirby Problem 4.43,'' including supplied literature triage dated 2026-08-17 and its orientation/invariant-notation warnings.
\item R. I. Baykur, R. C. Kirby, and D. Ruberman (eds.), \emph{K3}, AMS 295 (2026), Problem 4.43: \url{https://math.berkeley.edu/sites/default/files/surv-295-ruberman-watermarked-author-pdf.pdf}.
\item Jacob Rasmussen, ``Khovanov homology and the slice genus,'' \emph{Inventiones Mathematicae} 182(2) (2010), 419--447, DOI 10.1007/s00222-010-0275-6, source label [Ras10].
\item Ciprian Manolescu, Marco Marengon, Sucharit Sarkar, and Michael Willis, ``A generalization of Rasmussen's invariant, with applications to surfaces in some four-manifolds,'' \emph{Duke Mathematical Journal} 172(2) (2023), 231--311, source label [MMSW23]: \url{https://doi.org/10.1215/00127094-2022-0039}.
\item Michael Freedman, Robert Gompf, Scott Morrison, and Kevin Walker, ``Man and machine thinking about the smooth 4-dimensional Poincar\'e conjecture,'' \emph{Quantum Topology} 1(2) (2010), 171--208, DOI 10.4171/QT/5, source label [FGMW10].
\item P. B. Kronheimer and T. S. Mrowka, ``Gauge theory and Rasmussen's invariant,'' \emph{Journal of Topology} 6(3) (2013), 659--674, DOI 10.1112/jtopol/jtt008, source label [KM13b].
\item Sherry Gong, ``On the Kronheimer-Mrowka concordance invariant,'' \emph{Journal of Topology} 14(1) (2021), 1--28, DOI 10.1112/topo.12175, source label [Gon21].
\item Marco Mackaay, Paul Turner, and Pedro Vaz, ``A remark on Rasmussen's invariant of knots,'' \emph{Journal of Knot Theory and Its Ramifications} 16(3) (2007), 333--344, DOI 10.1142/S0218216507005312, source label [MTV07].
\item Robert Lipshitz and Sucharit Sarkar, ``A refinement of Rasmussen's S-invariant,'' \emph{Duke Mathematical Journal} 163(5) (2014), 923--952, DOI 10.1215/00127094-2644466, source label [LS14].
\end{itemize}
\end{document}
